\textbf{Время:} 01:15

\textbf{Тезис:} свойства источника исследуются через измеримые корреляции и проверяемую модель их связи с пространством и временем.

\textbf{Визуал:} пройти всю цепочку и связать каждый переход с его вопросом.

\textbf{Текст:} Детекторы дали треки, импульсы и идентификацию. Статистика пар добавила сведения, которых нет в отдельных спектрах. Бозонная интерференция связала относительный импульс с разделением точек испускания.

Так мы получили доступ к пространственно-временным масштабам выбранной области источника. Они помогают проверять сценарии расширения и работу транспортных моделей. Для вывода о динамике нужны те же пары, согласованные отборы и правильное описание взаимодействий.

Оставим короткий вопрос аудитории: в простом гауссовом примере размер увеличили вдвое. Как изменилась ширина импульсного сигнала? Дать ответить: уменьшилась вдвое. Затем предложить обсудить, какие дополнительные предположения потребовались бы для перехода от такого размера к длительности.
